\chapter{Programmablauf}
\label{chap:programmablauf}

% Inhalt laut Vorgabe:
%  - PAP oder Zustandsdiagramm o. Ä.
%  - ggf. Erläuterungen zum Bild

Die Steuerung von Modul~1 ist als Zustandsautomat umgesetzt. Die Methode
\texttt{schritt()} wird alle 0,1~s aufgerufen. Sie prüft zuerst die
Sicherheitsbedingungen und wertet danach den aktuellen Zustand aus.
Abbildung~\ref{fig:zustandsautomat} zeigt die Zustände und Übergänge.

\begin{figure}[htbp]
  \centering
  \begin{tikzpicture}[
      zustand/.style={draw, rectangle, rounded corners=4pt, minimum width=2.6cm,
                      minimum height=1.1cm, align=center, font=\small\bfseries},
      stopzustand/.style={zustand, draw=lamperot, very thick},
      uebergang/.style={-{Stealth[length=2.5mm]}, thick},
      bed/.style={font=\scriptsize, align=center},
      node distance=2.2cm and 3.2cm
    ]
    \node[zustand]                         (init)     {INIT};
    \node[zustand, right=of init]          (bereit)   {BEREIT};
    \node[zustand, right=of bereit]        (vorlauf)  {VORLAUF};
    \node[zustand, below=of vorlauf]       (laeuft)   {LÄUFT};
    \node[zustand, below=of bereit]        (nachlauf) {NACHLAUF};
    \node[stopzustand, below=of init]      (stop)     {STOP};

    \draw[uebergang] ($(init.west)+(-0.8,0)$) -- (init.west);

    \draw[uebergang] (init) -- node[bed, above] {Aktoren aus} (bereit);
    \draw[uebergang] (bereit) -- node[bed, above] {Silo nicht leer $\wedge$\\ M2 nicht voll}
                                node[bed, below] {/ Luft an} (vorlauf);
    \draw[uebergang] (vorlauf) -- node[bed, right] {nach 1~s\\ / Motor an} (laeuft);
    \draw[uebergang] (laeuft) -- node[bed, below] {Silo leer $\vee$\\ M2 voll} (nachlauf);
    \draw[uebergang] (nachlauf) -- node[bed, right] {nach 1~s\\ / Motor, Luft aus} (bereit);
    \draw[uebergang] (stop) -- node[bed, above, sloped] {kein Sicherheitsfehler} (bereit);

    \node[bed, below=0.9cm of stop, text width=5.2cm, draw=lamperot, dashed,
          rounded corners=3pt, inner sep=4pt] (alle)
      {aus jedem Zustand:\\ Not-Aus $\vee$ M2-Not-Aus $\vee$ Trichter voll\\ / Motor, Luft aus};
    \draw[uebergang, draw=lamperot] (alle) -- (stop);
  \end{tikzpicture}
  \caption{Zustandsautomat von Modul~1 (Delivery)}
  \label{fig:zustandsautomat}
\end{figure}

\section{Zustände}

\begin{table}[htbp]
  \centering
  \begin{tabularx}{\textwidth}{@{}l X@{}}
    \toprule
    \textbf{Zustand} & \textbf{Beschreibung} \\
    \midrule
    INIT     & Startzustand nach dem Einschalten. Alle Aktoren werden ausgeschaltet, danach folgt BEREIT. \\
    BEREIT   & Wartet, bis im Silo Material liegt und Modul~2 Material aufnehmen kann. Dann wird die Druckluft eingeschaltet. \\
    VORLAUF  & Die Druckluft läuft 1~s vor, damit der Schlauch frei ist, bevor das Förderband Mais nachschiebt. \\
    LÄUFT    & Förderband und Druckluft laufen. Ist das Silo leer oder Modul~2 voll, folgt der NACHLAUF. \\
    NACHLAUF & Förderband und Druckluft laufen 1~s nach, danach werden beide abgeschaltet. \\
    STOP     & Sicherheitszustand. Alle Aktoren sind aus, bis kein Sicherheitsfehler mehr anliegt. \\
    \bottomrule
  \end{tabularx}
  \caption{Zustände des Automaten}
  \label{tab:zustaende}
\end{table}

\section{Sicherheitsbedingungen}

Vor jeder Zustandsauswertung prüft der Automat drei Sicherheitsbedingungen:
den eigenen Not-Aus, den Not-Aus von Modul~2 und einen vollen Trichter. Trifft
eine davon zu, werden Motor und Druckluft im selben Schritt ausgeschaltet und
der Automat wechselt nach STOP, unabhängig vom aktuellen Zustand.

